%% APPENDIX D -- the false-alarm expectation, the control ensemble, and the
%% execution-block quality criterion.  Owner: r14-chance.  Carries
%% \label{app:falsealarm} (04_method.tex, 08_backmatter.tex) and
%% \label{sec:blockquality} (04_method.tex).
\section{The false-alarm expectation and the control ensemble}
\label{app:falsealarm}

This appendix gives the chance expectation for the unattributed crossings,
what it is conditioned on, and what the control ensemble calibrating it can
support.

\subsection{Conditions on the expectation}

\emph{One attribution set serves both sides of the comparison}, the set of
\S\ref{sec:technosearch}: the mask of Appendix~\ref{app:linemask} within
$\pm\MaskHalfKms$\,km\,s$^{-1}$ in the star's own rest frame, which holds no
sulphur monoxide. So the \SoNCrossWord{} crossings that fall on SO transitions
are unattributed crossings, and a control maximum on one of those transitions
is an unattributed exceedance: the same list on both sides, as it always was. A systemic velocity is measured for the stars of
\CeNVsys{} of the census windows and the others are taken at rest, which moves
the masked share by \CeVsysSensPct{} points.

The expectation is built from the control positions, the only empirical null
the data provide, and each Class~A window contributes two quantities
measured on itself. The fraction of its own \NCtrl{} control positions that
reaches the trigger is the probability that a position in that field does so
-- no cell count enters it and no form is assumed for the tail -- and
summed over the windows \CeExpTrig{} are expected to trigger by chance. How
many of those would be attributed is measured on the channels the control
maxima land on, which the retained spectra record: \CeNExcAttr{} of
\CeNExc{} exceedances fall inside the mask -- carried to the same star's rest
frame, through the same two steps and against the same transition list as a
cell at the stellar position -- and the attribution therefore
takes \CeMaskPct{} per cent of the expectation, where a maximum free to
land anywhere on its own grid would take \CeUnifPct{} per cent: the
\CeNFullMask{} narrow line-tuned windows lie wholly inside the mask and
their fields, which hold real emission, are where a control position most
often triggers, so a uniform landing would under-attribute them and leave
the unattributed expectation too high. The atmospheric screen of
Appendix~\ref{app:rfi} is then applied to these same positions, on the same
radius and in the same topocentric frame as to the crossings, and takes
\CeTlCtrlCut{} of the \CeTlCtrlSurv{} exceedances the mask leaves, a further
\CeTlPct{} per cent. That leaves \CeExp{} expected unattributed Class~A
crossings against \CeObs{} observed.

Both sides run over the \CeNWin{} Class~A windows of the primary census, so
the expectation is a sum over rows and nothing is carried at an assumed rate.
They are also read over the same field, and which field that is matters.
The $\sigma$ of Equation~\ref{eq:tstar} is the median absolute deviation over
the control positions, so $T_\star$ is a property of the annulus as much as of
the cell: one product supplies both sides here, the statistic read over the
whole primary beam with each block on its own antenna geometry, for the
expectation and for the crossing count alike. A field narrower than the
primary beam drops the integrations whose pointing falls outside it, and with
them on-source time and real signal, so a statistic read over part of a beam
is not the quantity \S\ref{sec:statistic} defines. The choice is not
immaterial and is reported rather than assumed small: read over such a field
on both sides, the same comparison gives \FrNarrowObs{} crossings against
\FrNarrowExp{} expected at $P=\FrNarrowProb$, where the whole beam gives
\CeObs{} against \CeExp{} at $P=\CeProbNull$; the bound below becomes
\FrNarrowAddLim{} further crossings against \CeAddLim{} on the same
resampling. The same comparison holds system by system, each system's
expectation formed from its own control positions: \MtNSysLowP{} of the
\MtNSysChance{} systems fall below one chance in twenty, against
\MtNSysLowPExp{} expected, and no system departs from its own controls by
more than chance allows.
A further \CeNWinTail{} searched Class~A windows lie outside the comparison
because the deposit does not tabulate them -- \VnTailNWinA{} in the
archival-tail blocks of \S\ref{sec:sample} and \VnCompNWinA{} toward a
companion resolved after the census closed -- together with the
\CeNTailCrossWord{} unattributed crossings toward \CeTailStar{} they hold,
which are reported beside it and enter neither side. No window holds more than one primary crossing, so
both sides count windows; the weaker cells behind those maxima are counted at
their own grain in \S\ref{sec:technosearch}.

The one assumption the searched data cannot test is that a star behaves like
a control position in its own field, and two samples measure it. In the census
the star carries \CeNear{} times as many above-threshold cells as a control
position of the same field, measured in the band immediately below the
trigger, $\CeNearBandLo\le T<\CeNearBandHi$, on \CeNearNStar{} stellar cells
over the \CeNearNWin{} windows that retain a control vector, with
\CeNearIntLo--\CeNearIntHi{} allowed at \CeCovPct{} per cent when the
execution blocks are resampled. The band ends at the trigger because the
ratio is not flat in level: read half-level by half-level it rises from
\CeNearRiseLo{} to \CeNearRiseHi, so a band further below the trigger would
not measure it where the comparison uses it. No cell under test enters it, a
band count being the difference of two cumulative counts, so an
above-threshold cell belonging to a crossing is counted in both terms and
cancels; at the trigger itself the ratio is \CeTrAt, but there it is formed
from the \CeStepNX{} cells that are the crossings under test. No step is
applied to the sum above: the expectation adopted is the control positions'
own, the smaller of the two and the more demanding of the count under test.
Carried in as a sensitivity test it raises the expectation to \CeExpNear{}
and the probability to \CeProbNear{} and lowers the bound to \CeAddLimNear,
the control positions' own rate being the one adopted.
What the comparison requires of that ratio is a statement in
the other direction: the \CeObs{} observed would stand above chance at the
\CeAlphaPct{} per cent level only if stars reached the trigger at \CeRCrit{}
times the rate of a control position or less, a deficit at the star of
\CeRCritPct{} per cent, and every ratio measured here is an excess instead.
The reserved blocks measure
the same quantity on blocks the search did not reach: their stars reach the
trigger in \StHoObs{} of \StHoNFine{} fine-channel windows where the control
positions of those same windows predict \StHoExpFine, a ratio of
\StHoRatio, and in the \StHoNCoarse{} coarse-channel reserved windows no
star reaches it against \StHoExpCoarse{} predicted. Every one of those
events lies in a reserved block that shares an observing session with a
searched one, and the \HsIndBlocks{} reserved blocks that share none hold
\HsIndObsWord{} trigger \HsIndObsNoun{} between them, so they neither
confirm the step nor contradict it (Appendix~\ref{sec:heldout}).

A count is compared with the distribution of counts chance delivers, and
not with a range on its own mean. A rate is drawn from \CeNBoot{} resamplings
-- each window's rate redrawn from its own \NCtrl{} probes, the attributed
share from its own \CeNExc{} exceedances, and the execution blocks resampled
because windows share them and are not independent trials -- and the count
scatters about that rate. A count of \CeObs{} or more then arises with
probability \CeProbNull{} at the control positions alone, and that is the
comparison adopted here; as a sensitivity test, applying the step gives
\CeProbStep. The same distribution bounds what the comparison could have
missed: added to each resampled control-position rate, \CeAddLim{} further
crossings still leave a count as low as \CeObs{} credible at \CeProbAtLim,
and one more does not, at \CeProbAboveLim{} against the \CeAlphaPct{} per
cent level. So \CeAddLim{} is the largest added population this comparison
could not have rejected, the uncertainty in \CeExp{} carried inside it by the
same resampling, and it is the compatibility statement of
\S\ref{sec:trigunif}: it counts crossings and not transmitters. Read at the grain of
the individual cell rather than the window, the same accounting gives
\CkExp{} expected against \CkNObs{} found, with \CkLo--\CkHi{} covering
\CeCovPct{} per cent of the same resampling.

\subsection{Execution-block quality}
\label{sec:blockquality}

The comparison with the control positions assumes each block's ensemble
behaves like noise, and one block makes that false: every one of its \NCtrl{}
probes exceeds the trigger, the smallest at \BqFailCtrlMin. A criterion for
the condition is therefore applied to every block and not to the one that
prompted it.

\emph{Two criteria are built from that statistic and they are not the same
test.} The statistic is the median of a window's own \NCtrl-position control
ensemble, which separates a field with a source in it from a field that is no
null at all. A \emph{window} is flagged when that median exceeds \BqFactor{}
the survey median of the same quantity, \DqThresh{} against \DqSurvMed; the
flag is a property of the window and cuts across attribution by design,
firing on \DqNWin{} of the \NWindows{} windows in \BkFlagNBlock{} blocks
toward \DqNStar{} stars (\DqStars), every one of which carries a crossing.
A \emph{block} is judged on the same median divided by the median over the
windows of its own class, since a maximum over a larger channel$\times$drift
grid sits higher; it fails where the median of that ratio over its windows
exceeds \BqFactor{} the scale of its class. The factor is the same in both.

\emph{Neither criterion was fixed in advance.} The condition was noticed
in one block and the criterion written afterwards, so it is post hoc, and
what makes it trustworthy is how little the position of the line matters. The
failing block sits at \BqFailQ{} and the next most extreme at \BqNextQ, a
factor \QsRatio{} apart; over \QsNStep{} thresholds spanning that interval
the same single block is selected every time, in number and in identity, so
the line may be drawn anywhere from \QsMarginLo{} to \QsMarginHi{} times
where it is. Outside that interval the answer does change, which is how one
can tell the rule discriminates: at \BqNextQ{} it selects \QsNAtNext{}
blocks, and at the median block, \QsQMed, \QsNAtMed{} of the
\BqNBlocks. Nor does the class treatment decide it: the same block is
selected with each class normalised on its own windows and with the survey
taken as one, and no block in the census is failed on one class
and passed on the other. The median block of the primary census sits at
\BqQMed, nothing else is close to the \BqNFail{} the criterion fails, and the
unscaled form of the rule fails the same \BqRawNFail{} block. Its
\BqFailNWin{} windows are all coarse-channel and
hold the survey's entire Class~B crossing population, so it carries
\EvNCrossFlag{} of the \EvNCrossRaw{} threshold crossings. \EvNRankFlag{}
of those outrank all \NCtrl{} control positions of their own window, and
no Class~A statistic here has ever contained it. Those \EvNCrossFlag{} appear
in the crossing ledger with their block named and enter no count
elsewhere.

\subsection{The control ensemble's resolution}

Over the \VerNTested{} census windows that retain a control vector the
stellar ranks are consistent with the one half exchangeability implies --
median \BkRankBalSurv{} with each star carrying equal weight,
\BkRankBalSurvLo{} to \BkRankBalSurvHi{} at \CeCovPct{} per cent -- which is
a weak statement and not a strong one, since a shift smaller than about
\VerKsPowerAca{} could not have been seen in a sample of this size, and the
reserved blocks agree once each star carries equal weight, \BkRankBalHo{}
(\BkRankBalHoLo--\BkRankBalHoHi). \emph{The rank is not a significance test,
and the reason is structural.}
Ranking the star against \NCtrl{} control positions of the same field
decides which windows deserve examination, but $1/(\NCtrl+1)=\RankFloorVal$
is the \emph{resolution} of the ensemble -- the smallest rank it can
express -- and not the probability that noise produced the window. A survey
of \NWindows{} windows would need a rank \BonfRankRatio{} times finer, and
correlation makes it coarser still: controls inside the synthesised beam are
not independent, so an annulus resolves no better than
\RankEffLo--\RankEffHi{} however many probes are drawn into it, with the
primary beam bounding how many a single pointing contains. That floor is the
design's central limitation: no single window here can be significant on the
rank alone, which is why confirmation requires independent recurrence.

\subsection{The trials budget}

The searched cells counted in \S\ref{sec:summary} are not independent
trials: channel smoothing and a drift grid stepped twice per channel
correlate neighbouring ones, which alone would leave \StGridFrac{} of them
independent, while fitting each window's effective count against its own
control ensemble returns \StNindFrac{} and those ensembles' maxima sit at
\CtrlMaxRatioMed{} of the prediction for the maximum of the full searched
count treated as independent and Gaussian. Nothing
turns on which is right, because the expectation above is summed from measured
rates and not from a cell count: by the cell-count route the number reaching
the trigger is \StExpModel{} against \CeExpTrig, \StModelAgreePct{} per
cent away. The on-star crossings do exceed a Gaussian tail over every
searched cell, which predicts \CellExpect{} stellar exceedances against the
\NCross{} found -- but \MkAttr{} are attributed by the set above,
\EvNCrossFlag{} come from the block of \S\ref{sec:blockquality},
\CeNTailCross{} fall outside the census windows and \CeTlObsCut{} is
atmospheric, which leaves \CeObs{}
against the \CeExpSurv{} the same control positions predict over both classes
together -- the \CeExp{} of the fine class and the \CeExpCoarseAdd{} the clean
coarse windows add to the expectation while holding \CeObsCoarseWord{} crossing
of their own, which is what keeps the comparison symmetric --
so no excess survives the attribution. The remainder is bounded and
the bounds are null: the measured stellar ranks agree at $\pv{\StarPseudoP}$
with the \NPseudo{} ranks each control takes against the other probes of its
own window, radial suppression is below \RingRhoBoundPct{} per cent in the
mean, and scrambling a window so nothing coherent survives along a drift
track carries \LNNWindows{} past the rank floor and changes no disposition.

%% The chance expectation here is summed from the measured control exceedance
%% rates and each window's own masked grid share (chance_v414.py, round 310),
%% so no cell count and no Gaussian tail enters it and nothing in it is
%% conditioned on the crossings; the cell-count route is a cross-check only,
%% and the rank-conditioned expectation (\ClustStageMean against
%% \ClustStageNObs) is a pipeline diagnostic that must not be reachable as
%% "the" expectation.  tab:ctrlsep, tab:excess, tab:arraysplit and
%% fig:ctrldiag are generated and nothing cites them; r15-deposit, 2026-10-08:
%% they used to be described as "in the data release", and there is no data
%% release, so they are simply unpublished diagnostics.
\subsection{The registration of the predicted cell}
\label{sec:regbudget}

The recurrence test of \S\ref{sec:etacrv} reads one channel at one drift
rate, so what it can conclude depends on how well that cell is registered
between two epochs. Every term of the budget is a frequency the reading
knows to a fraction of a channel: half a channel for the predicted cell,
half for the discovery channel it is referred to, \RgDriftWithinChMed{}
for the drift, \RgVsysChMax{} for the systemic velocity, which is removed at
one epoch and re-applied at the other, and \RgBaryChMax{} for the ephemeris,
a total of \RgBudgetChMed{} of a channel in every window and
\RgBudgetKmsMed\,km\,s$^{-1}$ in the median one. The frame correction is the
one term that is large in channels rather than small: it moves the predicted
cell by up to \ZtFrameChanMax{} channels, and by $\VnEcrDchan$ for
$\eta$~Crv. No reading at all is quoted for \RcUntStar{} \texttt{\RcUntEb},
whose frequency \RcUntNRep{} blocks cover but whose own window retains no
spectrum, so there is no discovery amplitude to transport to them. That is a
budget and not a
measurement, and the measurement is larger where the channels are narrow: the
stellar-frame centroid of a real recurring line sits a median
\RgCenMedKms\,km\,s$^{-1}$ from the predicted cell and \RgCenMaxKms{} at
worst, measured over \RgCenNBlock{} block readings of the line-attributed
recoveries. Each exclusion was therefore re-read over every cell within
\RgTolKms\,km\,s$^{-1}$ of the predicted one -- \RgNTrialMax{} cells at most
-- and formed from the worst of them, with the threshold raised to
\RgThrMax{} so that the wider reading keeps the single cell's own false-alarm
probability; it costs depth, the amplitude ruled out falling from a median
\RcFracMed{} to \RgFracMed{} of the discovery amplitude, and \RgNVac{} of the
\RcNExcl{} then rule out nothing, both of them already among the
\HjNPowerLo{} crossings \S\ref{sec:etacrv} reports as untested and not
excluded. Reading the same cell at every trial of the window's own drift grid
instead, which bounds an error in the carried drift whatever its size, leaves
the largest value anywhere \RcTMax{} over all \RcNRepeatMeas{} readings. The
reading has also to be pinned somewhere along the track, because a drifting
carrier occupies tens of channels and not one: each observation is de-drifted
from its own first integration, and pinning the cell anywhere across the span
the discovery track covers moves it by at most \AnSpanMaxKms\,km\,s$^{-1}$,
where the nearest above-trigger cell in any covering window lies
\AnNearKms\,km\,s$^{-1}$ away and the strongest reading over every anchor is
\AnTBest{} combined and \AnTSingle{} in a single block. So no disposition
rests on the carried drift, on the anchor or on the width of the cell.
Combining every covering block is what buys the depth that is quoted: the
nearest block alone rules out a median \TjFracNearMed{} of the discovery
amplitude, and \TjNNearWeak{} of the \RcNExcl{} rule out nothing on it. One
window is delivered twice by the archive, once with the control annulus on
that compact-array block's own antenna geometry and once on a nominal 12\,m
beam: the reading on the block's own geometry is the one published,
\LrSuppPub, where the 12\,m annulus gives \LrSuppOwn{} for the same cell,
\LrSuppPct{} per cent apart, \NbProvRmsPct{} per cent of that in the window
noise and the rest in the recovered amplitude. Part
of the power the single cell loses is the threshold and not the archive: a
statistic of five was set for a window of $10^{\TuCellsDecLo}$ to
$10^{\TuCellsDecHi}$ cells, and at
one predicted cell it is a far stricter demand, a false-alarm probability of
$\ZtOneCellP\times10^{\ZtOneCellExp}$ evaluated at the scatter measured in
\S\ref{sec:nullscatter}; holding the chance of a single false confirmation
across all \RcNRepeatMeas{} readings at \ZtFwPct{} per cent instead puts the
bar at \ZtTrialThr, and nothing would have been confirmed there either, the
largest single-block reading at the carried drift being \RcTMatched. A
carrier that keeps drifting across the gap is a different model and the
numbers belong to it: at the rates fitted here such a carrier is a median
\RgWalkChMed{} channels from the predicted cell by the nearest covering
block, and $\RgWalkChMax\times10^{\RgWalkChMaxExp}$ at \RgWalkStar.

\subsection{From the widened window's exceedances to the pairs}
\label{sec:pairacct}

Two counts of \S\ref{sec:etacrv} are counts of different objects, and
Table~\ref{tab:pairacct} takes one to the other. The \WdNExc{} above-trigger
readings are window readings, one for each pairing of an unattributed
crossing with a window that covers its frequency in another block, so a pair
whose two members are both unattributed crossings is read twice:
\PaNPairTwice{} of the readings are such a second reading, representing
\PaNPairRead{} distinct pairs, of which \PaNPairBeyond{} lie outside the
interval the pair test admits -- the median covering window's half-width --
and \PaNPairInside{} remain, which is the set Table~\ref{tab:pairs} is formed
over. The observatory's own line-of-sight acceleration is removed from both
fitted drifts before they are compared, at most \TjFrameMaxHzS\,Hz\,s$^{-1}$
here. Of the \TjNPower{} pairs whose rates the grid resolves, every one fails
a drift changing at a constant rate, by \TjSigMin{} standard deviations in
the mildest case; the offset
$\frac{1}{2}(\dot\nu_1+\dot\nu_2)\Delta t$ such a
trajectory predicts has no freedom left once both rates and the interval are
fixed, and the pair the searched grid itself separates is \EcGridStar's,
whose two crossings lie \EcGridDvKms\,km\,s$^{-1}$ apart in blocks
\EcGridDtH\,hours apart, a single carrier having had to drift at
\EcGridFactor{} times the grid's ceiling between them. Nothing measured here
bounds the curvature in between: the ceiling bounds the line-of-sight
component of a carrier's acceleration and not how fast that component can
turn over, an orbit directed across the plane of the sky sweeping it
\EcRateSupFac{} times faster than the circular orbit at the ceiling, while
the velocity excursion \emph{is} bounded, to within \EcVsupFac{} times the
circular amplitude over every eccentricity and orientation, which is why the
widened window of \S\ref{sec:etacrv} can be set against that excursion and a
reversal time cannot. The control row is a line
that does not drift: the \TjCtlLine{} line of \TjCtlStar{} is recovered at
one stellar-frame frequency \TjCtlDtD{} days apart and \TjCtlNSig{} standard
deviations from the prediction, its fitted drifts \TjCtlDriftStep{} of one
\TjCtlStepHzS\,Hz\,s$^{-1}$ grid step from zero and the predicted offset
below the comparison's own resolution over the same interval, so the
comparison resolves a mean drift to \TjCtlDriftResHzS\,Hz\,s$^{-1}$ and no
finer at any baseline. For the \TjNOpen{} pairs it cannot separate --
\TjNPowerOpen{} of them resolved and failed by the constant-acceleration
comparison -- nothing measurable tells one carrier from two.

$\beta$~Pictoris is the closest case, and it is the one the comparison
cannot dismiss. Its two crossings sit at \TjBpNuLo{}
and \TjBpNuHi\,GHz, \TjBpDvKms\,km\,s$^{-1}$ apart in the star's frame, in
blocks \TjBpDtH\,hours apart -- both inside the methanol $5_K\to4_K$ ladder,
which Appendix~\ref{app:linemask} excludes as their source -- with both
fitted drifts negative
($\TjBpDriftLo$ and $\TjBpDriftHi$\,Hz\,s$^{-1}$), so a single carrier would
have had to rise in frequency between the epochs and under constant
acceleration misses by \TjBpNSig{} standard deviations -- but the mean
acceleration it needs, \TjBpAccelReq\,m\,s$^{-2}$, is inside the
\WdAGrid\,m\,s$^{-2}$ ceiling. The pair fixes three numbers, the
acceleration at each epoch and the frequency offset between them, and a
circular orbit about a star of \EcMassBp\,M$_\odot$ \citep{Crifo1997} has
three free parameters: one orbit reproduces all three at once, at
\ObDayBp\,d in period, \ObIncBp$^\circ$ in inclination and \ObAuBp\,au in
radius, eccentric orbits fit up to $e$ of about \ObEccMaxBp{} before the
periastron acceleration leaves the ceiling, and the pair is therefore not
separable.

\begin{table}
%% ★ v4.26 (r25-captions): 5 printed lines -> the author's 3.  Naming
%% Table~\ref{tab:pairs} as what the readings are taken TO says in one clause
%% what two said: "the pairs the trajectory test is asked of" and "the counts
%% in the lower block are those of Table~\ref{tab:pairs}" are the same pairs.
\caption{The \WdNExc{} above-trigger readings of the widened recurrence
window, taken to the pairs of Table~\ref{tab:pairs}; each block of rows is
the one above with one effect removed.}
\label{tab:pairacct}
\footnotesize
%% GENERATED by pairacct_r23.py -- do not hand-edit.
\begin{tabular}{@{}p{0.76\columnwidth}r@{}}
\hline
Covering windows read, over the \PaNCovWinCross{} unattributed crossings a later block covers & \WdNCov \\
\ldots{} holding a statistic above the trigger, in \PaNRepWin{} windows of \PaNRepBlock{} blocks & \WdNExc \\
\hline
Distinct crossing pairs those readings represent & \PaNPairRead \\
\ldots{} readings that are the second reading of such a pair & \PaNPairTwice \\
Pairs separated by more than $\pm\WdHalfKms$\,km\,s$^{-1}$ in the star's frame, to \PaDvBeyondMax\,km\,s$^{-1}$ & \PaNPairBeyond \\
Pairs inside that interval & \PaNPairInside \\
\hline
Pairs with at least one unattributed member & \TjNCandPair \\
\ldots{} with a fitted drift and a block epoch for both members & \TjNPair \\
\hline
\ldots{} separated by the bounded-trajectory criterion & \TjNExcl \\
\ldots{} open & \TjNOpen \\
\hline
\end{tabular}

\end{table}

%% ★★ v4.23: THIS SUBSECTION AND ITS TABLE SIT AT THE END OF THE
%% APPENDIX FOR A REASON THAT IS NOT STYLE.  Placed next to the
%% registration budget, where it belongs by subject, its float took
%% the number Table~D1 and pushed the crossing ledger's sample to D2
%% and the recurrence positive controls to D3 -- and the positive
%% controls are the table Referee 2's first point is about, which he
%% and the response letter both call Table~D2.  Renumbering the
%% table a referee is reading about is a cost with no benefit, so the
%% new float goes last and takes the new number.  \S\ref{sec:edge}
%% is cited from \S4 and \S5.6 and the pointer is unaffected.
\subsection{The drift grid's edges}
\label{sec:edge}

The drift trials of a window run from $-\dot\nu_{\max}$ in equal steps of
half a native channel across the retained track, and the last trial is the
first at or beyond $+\dot\nu_{\max}$, so the interval searched is asymmetric
by up to one step and the ceilings themselves, computed window by window from
Equation~\ref{eq:accel}, reproduce it to $\DriftRelDev$ over all
\DriftRelNWin{} Class~A windows. Each window's grid is recovered from two
quantities that enter no common calculation -- the trial count recorded for
it and the channel width and retained span its own crossing record carries --
and the two agree on \EgNGrid{} of the \EgNDrift{} crossings that carry a
fitted drift, on every one of which the fitted rate is a trial of that grid
exactly; \EgNEdge{} lie on an outermost trial (Table~\ref{tab:edge}), toward
\EgEdgeStars, and are reported as lower bounds, while the remaining
\EgNNoGrid{} are in archival-tail windows for which no grid is recorded
anywhere.

\begin{table}
%% ★ v4.26 (r25-captions): 6 printed lines -> the author's 3.  DELETED after
%% checking: "a maximum at the boundary of the interval searched is a lower
%% bound" is \S\ref{sec:etacrv}'s own definition of \emph{edge-limited}, and
%% "Table~\ref{tab:pairs} gives them in the star's frame" is now a clause of
%% that same paragraph, which is where a reader meets both tables.
\caption{The crossings whose fitted drift lies on an outermost trial of
their window's own grid, with that grid. Rates are in the observatory
frame.}
\label{tab:edge}
\footnotesize
%% GENERATED by edge_v423.py -- do not hand-edit.
\begin{tabular}{@{}llr@{~}r@{~}c@{~}r@{~}l@{}}
\hline
Star & block & $\nu$ & $\dot\nu_{\rm obs}$ & trial & \multicolumn{2}{c}{$|a_{\rm los}|$} \\
 & & (GHz) & (Hz\,s$^{-1}$) & & \multicolumn{2}{c}{(m\,s$^{-2}$)} \\
\hline
BD+05~1668 & \texttt{X28a8} & 242.195 & $+3452$ & 2/2 & $\ge$ & 4.27 \\
HD~14055 & \texttt{X29ae} & 331.156 & $-3976$ & 0/135 & $\ge$ & 3.60 \\
HD~14055 & \texttt{X2c3bb} & 345.526 & $+4184$ & 141/141 & $\ge$ & 3.63 \\
HD~161868 & \texttt{X10ca8} & 345.634 & $+4195$ & 127/127 & $\ge$ & 3.64 \\
TRAPPIST$-$1 & \texttt{X259b} & 346.245 & $-4623$ & 0/137 & $\ge$ & 4.00 \\
\hline
\end{tabular}

\end{table}
